Los experimentos realizados confirmaron que la teoría asintótica predice bien el crecimiento general de los algoritmos, pero no refleja por sí sola el rendimiento real en la práctica, donde la gestión de memoria y la sobrecarga del código marcan una gran diferencia.

En el ordenamiento, el algoritmo más eficiente fue \textsc{QuickSort}, logrando los menores tiempos de ejecución y el uso más bajo de memoria al operar directamente sobre el arreglo original sin requerir estructuras adicionales. Por el contrario, el menos eficiente de forma general fue \textsc{PatienceSort}, que gastó mucha más memoria y tiempo debido a la sobrecarga de crear múltiples vectores dinámicos anidados en el heap. No obstante, este último mostró una notable ventaja adaptativa en arreglos ordenados de forma descendente y con dominios acotados ($D_1$), donde formó muy pocas pilas y alcanzó una rapidez lineal $O(n)$.

En la multiplicación de matrices, el método tradicional \textsc{Naive} superó con claridad a \textsc{Strassen} en matrices densas de hasta $1024 \times 1024$. Aunque la cota teórica favorece a \textsc{Strassen} con $\Theta(n^{2.807})$ frente al $\Theta(n^3)$ de \textsc{Naive}, el costo continuo de dividir la matriz en tantas llamadas recursivas y crear submatrices auxiliares anula el ahorro matemático. En conclusión, una solución con mejor cota teórica no siempre es la más rápida si su costo estructural supera al de una alternativa más simple.